\subsection{Summary \quad / \quad \textthai{บทสรุป}}
\begin{paracol}{2}
In this chapter, we progressed from passive observation to active control. The key insights gained include:
\begin{itemize}
    \item \textbf{Action-Value Mapping}: We used MDPs to formalize the goal of an agent as maximizing cumulative discounted rewards, framing physical stability as an optimization problem.
    \item \textbf{Bellman Optimality}: The Bellman Equation provides a recursive way to evaluate the "worth" of physical actions, allowing AI to plan multi-step trajectories.
    \item \textbf{DQN Success}: Deep Q-Networks successfully controlled high-dimensional systems (like the Rocket Landing) where classical linear control struggle with non-linear drift.
\end{itemize}

\switchcolumn

\begin{thai}
ในบทนี้ เราได้ก้าวข้ามจากการสังเกตการณ์ที่นิ่งเฉยไปสู่การควบคุมที่กระตือรือร้น ข้อมูลเชิงลึกหลักที่ได้รับ ได้แก่:
\begin{itemize}
    \item \textbf{การจับคู่ค่าการกระทำ}: เราใช้ MDP เพื่อกำหนดเป้าหมายของเอเยนต์เป็นการเพิ่มรางวัลสะสมแบบลดค่าให้สูงสุด โดยวางโครงสร้างเสถียรทางกายภาพเป็นปัญหาการหาค่าที่เหมาะสมที่สุด
    \item \textbf{ความเหมาะสมที่สุดของเบลล์แมน}: สมการเบลล์แมนเป็นวิธีการแบบเรียกซ้ำในการประเมิน "คุณค่า" ของการกระทำทางกายภาพ ช่วยให้ AI สามารถวางแผนวิถีแบบหลายขั้นตอนได้
    \item \textbf{ความสำเร็จของ DQN}: Deep Q-Networks ควบคุมระบบที่มีมิติสูง (เช่น การลงจอดของจรวด) ได้สำเร็จ ซึ่งการควบคุมเชิงเส้นแบบดั้งเดิมมักประสบปัญหาจากแรงบิดที่ไม่เป็นเชิงเส้น
\end{itemize}
\end{thai}
\end{paracol}
